% GENERATED FILE. Edit canonical lessons, then run npm run generate:books.
% canonical-source-hash: fnv1a64:01d4c0368e830930
% canonical-lessons: FR-C153-remarquer, FR-C153-remplacer, FR-C153-rencontrer, FR-C153-rougir, FR-C153-saigner

\chapter{To notice, To replace, To meet, To blush, To bleed}
\label{ch:fr-to-notice-to-replace-to-meet-to-blush-to-bleed}
\hlchaptermodality{153}

\begin{chapteropening}
\textbf{By the end of this chapter:} \emph{I can say to notice, to replace, to meet, to blush and to bleed.}

Complete the last lesson of chapter 153: I can say to notice, to replace, to meet, to blush and to bleed.
\end{chapteropening}

\section[remarquer]{remarquer — to notice}

\label{lesson:FR-C153-remarquer}



\begin{quote}
Before the new one: say the French for to move back, then the French for to lower.
\end{quote}

\subsection*{You'll want to know: remarquer}
\textbf{remarquer} — \textquotedblleft{}to notice\textquotedblright{}.

\begin{cousinweb}
\textbf{remarquer} — built from \emph{re-} and the verb for \textquotedblleft{}to mark\textquotedblright{}, whose noun came from a Germanic word for \textquotedblleft{}a sign\textquotedblright{}.
\end{cousinweb}

\begin{grammarlens}[title={What you've built}]
One more word for this chapter.
\end{grammarlens}

\subsection*{Your turn}
\begin{itemize}
\raggedright
  \item \emph{Say it:} \emph{remarquer}
  \item \emph{Say it:} \emph{remarquer}, once more
  \item \emph{Say it:} say \emph{remarquer}
  \item \emph{Recall:} say the French for to move back, then the French for to lower, then say \emph{remarquer} again
\end{itemize}

\begin{tcolorbox}[breakable,colback=teal!4,colframe=teal!35!black,title={Before you move on}]
What does remarquer mean? (\textquotedblleft{}To notice\textquotedblright{}.) Say it once more.
\end{tcolorbox}
\section[remplacer]{remplacer — to replace}

\label{lesson:FR-C153-remplacer}



\begin{quote}
Before the new one: say the French for to lower, then the French for to notice.
\end{quote}

\subsection*{You'll want to know: remplacer}
\textbf{remplacer} — \textquotedblleft{}to replace\textquotedblright{}.

\begin{cousinweb}
\textbf{remplacer} — built from \emph{re-}, \emph{en-} and the word for \textquotedblleft{}a place\textquotedblright{}, $\leftarrow$ Latin \emph{platea}, \textquotedblleft{}a broad street\textquotedblright{} (Lewis \& Short).
\end{cousinweb}

\begin{grammarlens}[title={What you've built}]
One more word for this chapter.
\end{grammarlens}

\subsection*{Your turn}
\begin{itemize}
\raggedright
  \item \emph{Say it:} \emph{remplacer}
  \item \emph{Say it:} \emph{remplacer}, once more
  \item \emph{Say it:} say \emph{remplacer}
  \item \emph{Recall:} say the French for to lower, then the French for to notice, then say \emph{remplacer} again
\end{itemize}

\begin{tcolorbox}[breakable,colback=teal!4,colframe=teal!35!black,title={Before you move on}]
What does remplacer mean? (\textquotedblleft{}To replace\textquotedblright{}.) Say it once more.
\end{tcolorbox}
\section[rencontrer]{rencontrer — to meet}

\label{lesson:FR-C153-rencontrer}



\begin{quote}
Before the new one: say the French for to notice, then the French for to replace.
\end{quote}

\subsection*{You'll want to know: rencontrer}
\textbf{rencontrer} — \textquotedblleft{}to meet\textquotedblright{}.

\begin{cousinweb}
\textbf{rencontrer} — built from \emph{re-} and Old French \emph{encontrer}, \textquotedblleft{}to meet\textquotedblright{}, from \emph{encontre}, \textquotedblleft{}against\textquotedblright{}, $\leftarrow$ Latin \emph{in contra} (Lewis \& Short for \emph{contra}).
\end{cousinweb}

\begin{grammarlens}[title={What you've built}]
One more word for this chapter.
\end{grammarlens}

\subsection*{Your turn}
\begin{itemize}
\raggedright
  \item \emph{Say it:} \emph{rencontrer}
  \item \emph{Say it:} \emph{rencontrer}, once more
  \item \emph{Say it:} say \emph{rencontrer}
  \item \emph{Recall:} say the French for to notice, then the French for to replace, then say \emph{rencontrer} again
\end{itemize}

\begin{tcolorbox}[breakable,colback=teal!4,colframe=teal!35!black,title={Before you move on}]
What does rencontrer mean? (\textquotedblleft{}To meet\textquotedblright{}.) Say it once more.
\end{tcolorbox}
\section[rougir]{rougir — to blush}

\label{lesson:FR-C153-rougir}



\begin{quote}
Before the new one: say the French for to replace, then the French for to meet.
\end{quote}

\subsection*{You'll want to know: rougir}
\textbf{rougir} — \textquotedblleft{}to blush\textquotedblright{}.

\begin{cousinweb}
\textbf{rougir} — built from \emph{rouge}, \textquotedblleft{}red\textquotedblright{}, with the verb ending \emph{-ir}: to go red. (EtymDB 2.1)
\end{cousinweb}

\begin{grammarlens}[title={What you've built}]
One more word for this chapter.
\end{grammarlens}

\subsection*{Your turn}
\begin{itemize}
\raggedright
  \item \emph{Say it:} \emph{rougir}
  \item \emph{Say it:} \emph{rougir}, once more
  \item \emph{Say it:} say \emph{rougir}
  \item \emph{Recall:} say the French for to replace, then the French for to meet, then say \emph{rougir} again
\end{itemize}

\begin{tcolorbox}[breakable,colback=teal!4,colframe=teal!35!black,title={Before you move on}]
What does rougir mean? (\textquotedblleft{}To blush\textquotedblright{}.) Say it once more.
\end{tcolorbox}
\section[saigner]{saigner — to bleed}

\label{lesson:FR-C153-saigner}



\begin{quote}
Before the new one: say the French for to meet, then the French for to blush.
\end{quote}

\subsection*{You'll want to know: saigner}
\textbf{saigner} — \textquotedblleft{}to bleed\textquotedblright{}.

\begin{cousinweb}
\textbf{saigner} $\leftarrow$ Latin \emph{sanguinare}, \textquotedblleft{}to bleed\textquotedblright{}, from \emph{sanguis}, the Latin of \emph{sang} (Lewis \& Short) — inherited.
\end{cousinweb}

\begin{grammarlens}[title={What you've built}]
One more word for this chapter.
\end{grammarlens}

\subsection*{Your turn}
\begin{itemize}
\raggedright
  \item \emph{Say it:} \emph{saigner}
  \item \emph{Say it:} \emph{saigner}, once more
  \item \emph{Say it:} say \emph{saigner}
  \item \emph{Recall:} say the French for to meet, then the French for to blush, then say \emph{saigner} again
\end{itemize}

\begin{tcolorbox}[breakable,colback=teal!4,colframe=teal!35!black,title={Before you move on}]
What does saigner mean? (\textquotedblleft{}To bleed\textquotedblright{}.) Say it once more.
\end{tcolorbox}
